\documentclass[11pt,a4paper]{article}
\usepackage{fontspec}
\setmainfont{Noto Serif}
\setsansfont{Noto Sans}
\usepackage{babel}
\babelprovide[main,import]{russian}
\usepackage{amsmath,amssymb,mathtools,booktabs,geometry,fancyhdr,microtype,xcolor,hyperref}
\geometry{left=21mm,right=21mm,top=21mm,bottom=21mm}
\definecolor{ink}{HTML}{173D57}
\hypersetup{colorlinks=true,urlcolor=ink,linkcolor=ink,pdftitle={BGS-F. Точность квантования и гейзенберговские следы}}
\pagestyle{fancy}\fancyhf{}\lhead{\small BGS-F}\rhead{\small Точность квантования и гейзенберговские следы}\cfoot{\thepage}
\setlength{\headheight}{15pt}\setlength{\parindent}{1em}\setlength{\parskip}{2.5pt}
\setlength{\emergencystretch}{3em}\tolerance=3000
\newcommand{\Tr}{\operatorname{Tr}}
\newcommand{\EE}{\mathbb E}
\newcommand{\R}{\mathbb R}
\newcommand{\Z}{\mathbb Z}
\newcommand{\C}{\mathbb C}
\newcommand{\si}{\mathrm{Sine}_2}
\newcommand{\eps}{\varepsilon}
\newcommand{\sinc}{\operatorname{sinc}}
\newcommand{\sect}[1]{\section*{\sffamily\color{ink}#1}}
\newcommand{\subt}[1]{\par\medskip\noindent\textbf{#1.}\quad}
\newcommand{\T}{\mathbb T}
\newcommand{\dist}{\operatorname{dist}}
\newcommand{\Ad}{\operatorname{Ad}}
\newcommand{\zp}{\mathbb Z_7}
\newcommand{\QQ}{\mathbb Q}
\newcommand{\kk}{\overline{\mathbb Q}}
\newcommand{\ord}{\operatorname{ord}}
\newcommand{\diag}{\operatorname{diag}}
\begin{document}
\noindent{\small\sffamily\color{ink} МАТЕМАТИЧЕСКАЯ ФИЗИКА\hfill 22 НОЯБРЯ 2026}
\begin{center}{\LARGE\sffamily Точность квантования и\\[3pt]гейзенберговские следы}\\[6pt]
{\large Устойчивость локального закона и резкий порог в норме}
\end{center}
Рассматриваем исходную задачу BGS-F для \(U_N=\exp(iN10^{-80} A_N)M_N\). Смысл локального предела и обозначение \((P)\) взяты из постановки. Доказательства ниже разделяют два класса количественных входных оценок, которые могут использоваться при замене матрицы аппроксимацией.

\subt{Перенос закона}
Для двух последовательностей унитарных матриц \(U_N,V_N\) условие \(\|U_N-V_N\|_{S_1}=o(1)\) даёт одинаковые локальные слабые пределы с общим равномерным сдвигом. Для исходного парного теста обозначим через \(\gamma_U\) меру разностей с диагональю, а через \(g_*\) автокорреляцию функции \(b_*\). При \(N\|U_N-V_N\|_{\rm op}\) достаточно малом справедлива относительная граница
\[
    |\gamma_V(g_*)-\gamma_U(g_*)|
       \le C_*N\|U_N-V_N\|_{\rm op}\,\gamma_U(g_*).
\]
Значит, ошибка \(o(1/N)\) переносит в обе стороны и сам тест \((P)\), без отдельного априорного предположения о втором моменте.

\subt{Пример на пороге}
Для любой исходной последовательности \(U_N\), в частности для фиксированного \(10^{-80}\), построены соседние вспомогательные \(V_N\) в том же собственном базисе, с тем же определителем и \(\|V_N-U_N\|_{\rm op}=O(1/N)\). Спектр \(V_N\) прост, угловые зазоры между различными фазами не меньше \(3/(4N^3)\) в радианах, а все нетривиальные сбалансированные целые мономы по собственным числам трансцендентны. Можно выбрать \(V_N\) так, чтобы абсолютные ошибки коротких следов стремились к нулю равномерно на любом заданном окне \(K_N=o(\sqrt N)\). Для \(K_N=o(N^{2/5})\) также стремится к нулю сумма квадратов этих ошибок. При этом на исходном весе \(w\) для вспомогательной последовательности
\[
                     \liminf_N\Delta_{2,N}(V_N;b_*)\ge101/32 .
\]
Другая проверка использует ограниченную непрерывную статистику локального окна. Она отделяет \(V_N\) от \(\si\) независимо от поведения вторых моментов.

\subt{Точное совпадение следов до половины размера}
Для \(N=2q+1\) приведены два простых спектра с определителем 1 и точно одинаковыми следами для \(1\le m<q\). Их предельные значения парного теста различаются на \(\tfrac34[w(1/2)-w(1)]>0\). Их локальные законы тоже различны.

В примерах меняется квантовый оператор. Для исходной экспоненциальной формулы при фиксированной силе предел Sine\(_2\) и условие \((P)\) остаются открытыми.

\newpage
\sect{1. Операторное расстояние и перенос точечного процесса}
Пусть \(U,V\in U(N)\). Запишем \(VU^*=e^{iH}\), где \(H\) эрмитова и её собственные числа лежат в \([-\pi,\pi]\). Рассмотрим \(W(s)=e^{isH}U\), \(0\le s\le1\). Углы здесь измеряются в радианах. По теореме Реллиха можно локально выбрать аналитические собственные ветви, взяв эрмитов логарифм с разрезом вне спектра в каждой окрестности параметра. На концах этих окрестностей ветви сопоставляются по значениям. Получаем подъёмы \(\varphi_j(s)\) собственных углов и соответствующие ортонормированные векторы \(v_j(s)\) на каждом таком участке. Дифференцирование унитарной задачи даёт
\[
      \dot\varphi_j(s)=\langle v_j(s),H v_j(s)\rangle .
\]
Формула работает и при совпадении фаз на аналитической ветви. Следовательно,
\[
 \sum_j|\dot\varphi_j(s)|\le\Tr|H|,qquad
 \max_j|\dot\varphi_j(s)|\le\|H\|_{\rm op}.
\]
При интегрировании это даёт единое сопоставление фаз \(U,V\), одновременно для суммы и максимума угловых смещений. Неравенство \(|x|\le(\pi/2)|e^{ix}-1|\) при \(|x|\le\pi\) переводит оценки в матричные нормы. После деления на \(2\pi\), в единицах поворота имеем
\begin{equation}
  \sum_j|\theta_j'-\theta_j|\le\tfrac14\|U-V\|_{S_1},
  \qquad
  \max_j|\theta_j'-\theta_j|\le\tfrac14\|U-V\|_{\rm op} .
                                                        \tag{1}
\end{equation}
Здесь используются выбранные подъёмы углов, поэтому периодизация исключает проблему разреза окружности.

\subt{Общий сдвиг}
Пусть \(f\in C_c^1(\R)\). При общем равномерном \(\alpha\) для двух спектров интегрируем разность соответствующих линейных статистик по \(\alpha\). Сумма по целым периодам переводит интеграл в интеграл по всей прямой. Для каждого слагаемого производная контролирует сдвиг на \(N|\theta_j'-\theta_j|\), а замена переменной в интеграле даёт множитель \(1/N\). Поэтому
\begin{equation}
 \EE_\alpha|\Xi_{V,\alpha}(f)-\Xi_{U,\alpha}(f)|
       \le \tfrac14\|f'\|_{L^1}\|U-V\|_{S_1}.            \tag{2}
\end{equation}
При \(f\ge0\) та же граница контролирует разность функционалов Лапласа, поскольку \(x\mapsto e^{-x}\) липшицева на \([0,\infty)\). Таким образом, условие \(\|U_N-V_N\|_{S_1}=o(1)\) переносит локальный слабый закон. В частности, достаточно \(\|U_N-V_N\|_{\rm op}=o(1/N)\). Для неунитарной аппроксимации \(A_N\) с \(\|A_N-U_N\|_{\rm op}=o(1/N)\) унитарный полярный множитель отличается от \(U_N\) по норме не больше \(2\|A_N-U_N\|_{\rm op}\), поскольку все сингулярные числа отличаются от 1 не больше, чем на эту исходную ошибку без множителя 2.

\newpage
\sect{2. Устойчивость именно парного теста}
Для списка фаз \(\theta_j\) с кратностями введём положительную меру
\[
  \gamma_U=\frac1N\sum_{j,k=1}^N\sum_{l\in\Z}
             \delta_{N(\theta_j-\theta_k+l)} .
\]
Положим \(g_*(x)=\int b_*(y)b_*(y+x)\,dy\), где \(b_*,w\) точно как в постановке \((P)\). Тогда \(g_*\) --- функция Шварца, строго положительная на всей прямой. Действительно, обратное преобразование Фурье гладкой компактно поддержанной ненулевой \(\phi\) есть целая функция, чьи вещественные нули изолированы. В произведении двух сдвигов \(b_*\) интегранд положителен почти всюду. Суммирование Пуассона и \(\widehat{g_*}=w\) дают
\[
       2\Delta_{2,N}(U;b_*)
               =\gamma_U(g_*)-\EE_{\rm CUE}\gamma_{\rm CUE}(g_*).
\]
Диагональный вклад и нулевой Фурье-член совпадают и вычитаются. Последнее слагаемое равномерно ограничено, поскольку \(w\) ограничен и поддержан на \([-2,2]\), а средний квадрат ненулевого следа CUE равен \(\min(|m|,N)\).

\subt{Верхняя граница хвостовых масс из одного теста}
При \(h(x)=(1-|x|)_+\) и \(Y_a=\Xi_{U,\alpha}([a,a+1])\) стационарность по общему сдвигу даёт
\[
  (\gamma_U*h)(t)=\EE Y_0Y_t\le\EE Y_0^2=(\gamma_U*h)(0).
\]
На интервале длины 1 вокруг \(t\) функция \(h(t-\cdot)\) не меньше \(1/2\). Кроме того, \(h\le Cg_*\) на её компактной опоре. Значит, масса \(\gamma_U\) каждого единичного интервала не превосходит \(C\gamma_U(g_*)\), с одной и той же константой по \(N,t\). Для
\[
                 G(x)=\sup_{|y-x|\le1}|g_*'(y)|
\]
ряд супремумов \(G\) по единичным интервалам конечен. Поэтому \(\gamma_U(G)\le C_*\gamma_U(g_*)\).

\subt{Относительная ошибка}
По (1) каждая сопоставленная развёрнутая разность фаз сдвигается не больше чем на \(N\|U-V\|_{\rm op}/2\). Если это число не превосходит 1, теорема о среднем и суммирование по всем парам дают
\[
 |\gamma_V(g_*)-\gamma_U(g_*)|
           \le \tfrac12 N\|U-V\|_{\rm op}\,\gamma_U(G)
           \le C_*N\|U-V\|_{\rm op}\,\gamma_U(g_*).
\]
Отсюда следует двусторонняя устойчивость \((P)\) при \(o(1/N)\) в операторной норме. Если условие выполнено хотя бы для одной из последовательностей, равномерная ограниченность её \(\gamma(g_*)\) следует из самого теста, и абсолютная разность стремится к нулю. Подобная оценка работает для другого функционала пары с производной, у которой суммируются единичные супремумы, как только \(\gamma_U(g_*)\) ограничена.

\newpage
\sect{3. Округление с фиксированным определителем}
Углы исходной матрицы в радианах обозначим \(\varphi_j\). Для больших \(N\) возьмём
\[
           q=\lfloor N/32768\rfloor,\qquad h=2\pi/q,
           \qquad \varphi_j=a_j+h t_j,quad 0\le t_j<1,
\]
где \(a_j\) --- левый соседний узел сетки \(h\Z\) в выбранном подъёме. Пусть \(K=K_N\) --- заданное число коротких следов. До округления задаём линейные условия на \(p=(p_1,\ldots,p_N)\in[0,1]^N\). Для каждого \(1\le m\le K\) сохраняем комплексное число
\[
   \sum_j\left((1-p_j)e^{im a_j}+p_j e^{im(a_j+h)}\right)
\]
с тем же значением, что при \(p=t\). Дополнительно сохраняем \(\sum_jp_j=\sum_jt_j\). Последняя строка сохранит математическое ожидание полного поднятого угла. Система выполнима при \(p=t\).

В вершине этого многогранника не больше \(2K+1\) дробных координат. Если их больше, существует ненулевой допустимый в обе стороны вектор в ядре \(2K+1\) вещественных ограничений, поддержанный на этих координатах. Это противоречит экстремальности. Поэтому выбираем такую вершину и независимо заменяем оставшиеся дробные координаты случайными \(B_j\in\{0,1\}\) с \(\EE B_j=p_j\). Целые координаты оставляем. Пусть \(r\le2K+1\) --- число дробных переменных, \(\psi_j=a_j+hB_j\). Все \(\psi_j\) лежат на общей сетке.

\subt{Ожидание и дисперсия}
Обозначим \(Z_m=\sum_j e^{im\psi_j}\), \(p_m(U)=\Tr U^m\) и \(S=\sum_j(\psi_j-\varphi_j)\). Суммарные смещения в ожидании равны нулю, поскольку добавлено последнее линейное условие. При этом
\[
  |\EE Z_m-p_m(U)|\le\tfrac12 Nm^2h^2,\qquad
  \operatorname{Var}(Z_m)\le r m^2h^2,qquad
  \EE S^2\le r h^2.
\]
Первая оценка следует из линейной интерполяции в исходных точках \(\varphi_j\), где линейный член в разложении экспоненты сокращается. При вычислении дисперсии участвуют лишь оставшиеся случайные координаты. Их разности по амплитуде не больше \(mh\).

Уберём из каждого \(\psi_j\) общий угол \(S/N\). Тогда определитель совпадает с \(\det U\) при каждом исходе. Модуль изменения \(m\)-го следа при общей поправке не больше \(m|S|\), поскольку \(|Z_m|\le N\). Поэтому
\begin{equation}
 \EE\sum_{m=1}^{K}|e^{-imS/N} Z_m-p_m(U)|^2
   \le 4 r h^2\sum_{m=1}^{K}m^2+\frac12N^2h^4\sum_{m=1}^{K}m^4 .
                                                        \tag{3}
\end{equation}
Это \(O(K^4/N^2+K^5/N^2)\). Для одного подходящего исхода эта граница сохраняется с абсолютным множителем 4. В следующем разделе тот же исход выбирается с дополнительным равномерным контролем.

\newpage
\sect{4. Все короткие следы и арифметика спектра}
\subt{Равномерный контроль}
Для случайной вещественной величины \(Y\) с размахом не больше \(D\) и независимой копии \(Y'\) неравенство Йенсена даёт \(\EE e^{u(Y-\EE Y)}\le\EE e^{u(Y-Y')}\le\exp(u^2 D^2/2)\), используя симметрию разности и \(\cosh x\le e^{x^2/2}\). Суммирование независимых величин и оптимизация по \(u\) дают двусторонний субгауссов хвост. Применяем его отдельно к действительной и мнимой частям \(Z_m-\EE Z_m\), а также к \(S\). Размахи равны сверху \(mh\) и \(h\). Объединяя \(2K+1\) событий, с вероятностью не меньше \(3/4\) одновременно имеем
\[
 |Z_m-\EE Z_m|\le C m h\sqrt{r\log(2K+2)},\qquad
 |S|\le C h\sqrt{r\log(2K+2)}.
\]
Событие с квадратной оценкой (3), умноженной на 4, также имеет вероятность не меньше \(3/4\). Выбираем исход из пересечения. Таким образом, после общей поправки абсолютная максимальная ошибка не больше
\[
              O\!\left(\frac{K^2}{N}
                         +\frac{K^{3/2}\sqrt{\log(2K+2)}}{N}\right).
\]
Она стремится к нулю при \(K=o(\sqrt N)\). При \(K=o(N^{2/5})\) правая часть (3) стремится к нулю. Выбор исхода нужен только для существования детерминированного семейства, исходные спектральные пределы случайны лишь по равномерному общему сдвигу.

\subt{Простота, разделение и трансцендентность}
Пусть \(p_j\) теперь обозначает \(j\)-е простое число. К каждому поднятому углу добавим малую величину
\[
          \delta_j=\frac{j}{N^3}
                   +\frac{\sqrt{p_j}}{4N^3p_N}
\]
и вычтем их общий средний угол \(\bar\delta\). Определитель сохраняется. Внутри каждого узла сетки различные углы отличаются как минимум на \(3/(4N^3)\), потому что рациональная часть меняется на \(1/N^3\), а общая ширина добавки с корнем не превосходит \(1/(4N^3)\). Между различными узлами остаётся расстояние не меньше \(h-2/N^2\), включая шов окружности. Для ненулевого целого вектора \(n\) с \(\sum_j n_j=0\) равенство \(\sum_j n_j\delta_j=0\) невозможно. Действительно, \(1,\sqrt{p_1},\ldots,\sqrt{p_N}\) линейно независимы над \(\QQ\) по независимым сменам знаков в мультиквадратичном поле. Соответствующий моном в собственных числах равен корню единицы, умноженному на экспоненту ненулевого алгебраического числа. Он трансцендентен по Эрмиту--Линдеману.

Вспомогательную матрицу \(V_N\) задаём этими углами в исходном собственном базисе. Джиттер меняет \(m\)-й след не больше \(2m/N\), поэтому оценки короткого окна сохраняются. Операторное расстояние ограничено \(h(1+r/N)+2/N^2\). Здесь \(|S|\le rh\) даже без вероятностной оценки, поскольку ошибки интегральных переменных равны нулю после вычитания их ожиданий. При \(K=o(\sqrt N)\) получаем \(O(1/N)\).

\newpage
\sect{5. Нарушение на заданном весе и слабый тест}
Для \(m=bq\) с \(1\le b\le8192\) вклад всех узлов сетки после целой степени совпадает по фазе. Общие поправки дают общий множитель модуля 1. Остаётся сумма \(e^{im\delta_j}\), откуда
\[
         \frac{|\Tr V_N^{bq}|}{N}\ge1-\frac{1}{2N},
       \qquad 0<\frac{bq}{N}\le\frac14 .
\]
Здесь используется \(\max_j\delta_j\le2/N^2\). Для \(0\le t\le1/4\) в свёртке \((\phi*\phi)(t)\) интегрируем по \([-1/4,1/4]\). Оба аргумента по модулю не больше \(1/2\). На этом интервале \(\phi\ge e^{-4/3}>1/4\), а \(4I^4\le4(2/e)^4<2\). Значит, \(w(t)\ge1/2048\). Для всех \(t\) имеем \(w(t)\le e^{-2}/(4I^2)<9/16\), поскольку \(I\ge e^{-4/3}\) и \(e^{1/3}<3/2\). Поэтому точная CUE-часть парного теста не больше \((9/16)\cdot(3/2)=27/32\). Оставляя только \(8192\) степеней \(bq\) в положительной части, получаем
\[
        \liminf_N\Delta_{2,N}(V_N;b_*)\ge
                          8192/2048-27/32=101/32 .
\]
Здесь использован именно заданный в задаче вес, с полной вычитаемой рампой \(\min(m,N)\).

\subt{Ограниченное наблюдение}
Пусть \(f\) непрерывна, поддержана на \([-1,1]\), равна 1 на \([-1/2,1/2]\) и лежит между 0 и 1. Для \(F(\xi)=\min(1,\xi(f))\) получаем
\[
                 \limsup_N\EE_\alpha F(\Xi_{V_N,\alpha})\le1/16384.
\]
Действительно, все точки до развёртки лежат на расстоянии \(O(N^{-2})\) от общих \(q\) узлов в единицах поворота. При равномерном сдвиге вероятность присутствия хотя бы одной точки в \([-1,1]\) не больше \(q(2/N+O(N^{-2}))\). Для \(\si\) и счётчика \(X\) на \([-1/2,1/2]\) имеем \(F\ge\mathbf1_{X\ge1}\ge X-\binom{X}{2}\). Так как \(\EE X=1\) и двухточечная интенсивность не превосходит 1, то \(\EE_{\si}F\ge1/2\). Значит, слабый закон тоже исключён для построенной последовательности.

\subt{Одношаговые и коротковременные оценки}
Положим \(P_U=(|U_{xy}|^2)_{xy}\). Оно есть сжатие \(U\otimes\bar U\) изометрией \(e_x\mapsto e_x\otimes e_x\), поэтому \(\|P_V-P_U\|_{\rm op}\le2\|V-U\|_{\rm op}\). Тем же способом сохраняется асимптотическое убывание нормы обратного канала \(U\circ U^*\). Верхняя граница \(O(N^{-1/2})\) для элементов и одноступенчатая оценка Егорова с остатком \(O(N^{-1})\) сохраняются по треугольнику. Кроме того, на окне \(m=o(N)\) нормированные следы отличаются на \(o(1)\) по телескопическому неравенству \(|\Tr V^m-\Tr U^m|/N\le |m|\|V-U\|_{\rm op}\). Простота даёт точное бесконечное среднее квадрата следа \(N\) для каждой матрицы \(V_N\).

\newpage
\sect{6. Точное совпадение почти половины следов}
Рассмотрим \(N=2q+1\) и параметр \(a\in(0,1/2]\). Возьмём две сетки и одну дополнительную фазу
\[
 \left\{\frac{j-a/2}{q},\frac{j+a/2}{q}\colon0\le j<q\right\}\cup\{0\}
                  \quad\text{в }\R/\Z .
\]
Это список собственных фаз диагональной унитарной матрицы \(C_{q,a}\), в единицах полного оборота. У обеих сеток произведения собственных чисел равны \((-1)^{q-1}e^{\mp\pi i a}\). Поэтому полный определитель равен 1. Фазы просты, а наименьший зазор равен \(a/(2q)\) из-за добавленной фазы. Суммирование геометрических прогрессий даёт при любом \(m\ge1\)
\[
      \Tr C_{q,a}^m=1+2q\,\mathbf1_{q\mid m}\cos(\pi a m/q).
\]
Значит, при разных \(a\) точно совпадают все следы \(1,\ldots,q-1\). Для предельной разности парного теста у \(a=1/6\) и \(1/2\) остаются только индексы \(m=q,2q,3q,4q\). Последний вклад стремится к нулю, потому что \(w(2)=0\). После деления на \(N^2\) кросс-члены с единицей исчезают. Получаем
\[
  \lim_{q\to\infty}\!
  \left[\Delta_{2,2q+1}(C_{q,1/6};b_*)
              -\Delta_{2,2q+1}(C_{q,1/2};b_*)\right]
       =\tfrac34\bigl(w(1/2)-w(1)\bigr)>0.
\]
Строгое неравенство следует из представления \(\phi\) интегралом индикаторов центрированных интервалов. Длины пересечений этих интервалов при относительном сдвиге не возрастают. Для открытого множества пар радиусов, близких к \(0.7\), длина строго уменьшается между сдвигами \(1/2\) и 1. Уровни \(\phi\), соответствующие этим радиусам, образуют множество положительной меры.

\subt{Локальные законы}
После развёртки две сетки имеют период \(N/q\to2\), а два смещения относительно центра стремятся к \(\pm a\). Равномерный общий сдвиг на окружности длины \(N\) после взятия остатка по этому периоду равномерен. Дополнительная фаза попадает в любое фиксированное окно с вероятностью порядка \(1/N\) и исчезает из слабого предела. При \(a=1/2\) предельные интервалы равны 1, при \(a=1/6\) они чередуются как \(1/3,5/3\). Таким образом, общие фазовые данные через \(q-1\) и простой спектр с зазором порядка \(1/N\) совместимы с разными локальными законами и значениями заданного парного теста.

\newpage
\sect{7. Связь с ранее доказанными оценками}
\subt{Что добавлено к округлению проекта}
Два предыдущих результата в репозитории описаны в \texttt{balanced\_rounding\_research\_20261027.md} и \texttt{microscopic\_rounding\_control.md}. Первый сохраняет малое абсолютное окно следов детерминированным балансированием при более грубой норме ошибки. Второй использует независимое округление и одновременно контролирует норму \(O(1/N)\), квадратную ошибку на окне \(o(N^{1/3})\) и отдельную гладкую пробу в стабильной полосе. Здесь балансирование сокращает число случайных координат до \(2K+1\), даёт квадратную оценку на \(o(N^{2/5})\) и равномерную абсолютную на \(o(\sqrt N)\). Сохранён определитель. Нарушение доказано для исходного \(w\) вместе с независимым ограниченным наблюдением. Полные тексты использованных замечаний включены в папку с указанием исходной ревизии.

\subt{Зачем нужна точная одношаговая информация}
Фиксированное классическое отображение, сжатия одношаговых амплитуд, абсолютные короткие следы, простота, отсутствие точных алгебраических соотношений и грубая нижняя граница зазоров совместимы с сильным нарушением \((P)\). Конструкция демонстрирует это одним семейством \(V_N\), построенным около фактического \(U_N\). Для самого \(U_N\) сила \(10^{-80}\) и определение \(A_N\) сохранены в постановке. Для доказательства \((P)\) требуется модельная оценка когерентного суммирования при \(m\asymp N\). Для положительного слабого предела дополнительно нужны все локальные функционалы, например посредством полной факториальной иерархии.

\subt{Воспроизведение}
Команда \texttt{python3 resolution\_precision/check.py --verify} пересчитывает рациональные границы на вес, слагаемые в обоих парных тестах, арифметику сеток и суммы в оценке (3). Скрипт \texttt{resolution\_precision/check\_matrix.py} на умеренной размерности независимо строит вершину многогранника линейным программированием, выбирает округление условным математическим ожиданием и проверяет численно сохранение определителя и одношаговой оценки. Эти вычисления проверяют реализации и константы. Асимптотические утверждения доказаны выше без численных предположений.

\subt{Ссылки}
Теорема Реллиха для аналитических эрмитовых семейств приведена в обзоре о возмущениях многочленов и линейных операторов, теорема 2.1. \href{https://www.mat.univie.ac.at/~armin/publ/handbook_paper.pdf}{Текст обзора.}

Прежние результаты об округлении и статус самой BGS-F приведены в \href{https://github.com/ChurilovMaksim/BGS/tree/e2757e4cc5ef996e08190f88c93c710377c10c22}{зафиксированной ревизии BGS}. См. также \texttt{research/problem.md} в папке для полной постановки.

\end{document}
